% Generated draft: requires mathematical and source review.
\documentclass[11pt]{article}
\usepackage[margin=25mm]{geometry}
\usepackage{fontspec}
\setmainfont{Latin Modern Roman}
\usepackage{amsmath,amssymb,mathtools,unicode-math}
\setmathfont{Latin Modern Math}
\usepackage{enumitem,needspace,xurl,hyperref}
\setlength{\parindent}{0pt}
\setlength{\parskip}{5pt}
\setlength{\emergencystretch}{3em}
\newcommand{\N}{\mathbb{N}}
\newcommand{\Z}{\mathbb{Z}}
\newcommand{\Q}{\mathbb{Q}}
\newcommand{\R}{\mathbb{R}}
\newcommand{\C}{\mathbb{C}}
\newcommand{\F}{\mathbb{F}}
\newcommand{\A}{\mathbb{A}}
\newcommand{\PP}{\mathbb{P}}
\newcommand{\E}{\mathbb{E}}
\newcommand{\eps}{\varepsilon}
\newcommand{\dd}{\,\mathrm{d}}
\begin{document}
\section*{UnsolvedMath ID: 30002847}
\textbf{Source problem code:} catalogue entry
\subsection{Definitions and mathematical statement}
The supplied question concerns the Ihara zeta function of a connected regular graph. The record does not explicitly specify finiteness, directedness, or whether loops and multiple edges are allowed. The following definitions give a conditional formulation under the conventional finite, undirected, simple graph interpretation; these conventions require confirmation against the source.

Let \(G=(V,E)\) be a finite, connected, undirected, simple \(k\)-regular graph, where \(k\in\Z_{\geq 0}\), meaning that every vertex has degree \(k\). Put \(n=|V|\) and \(m=|E|\). Let \(A\in\{0,1\}^{n\times n}\) be its adjacency matrix, and let \(I_n\) be the identity matrix. Replace each edge by its two orientations, and denote the reverse of an oriented edge \(e\) by \(\bar e\).

A closed oriented edge walk \(C=(e_1,\ldots,e_\ell)\) is cyclically nonbacktracking if successive edges are incident in the direction of travel and \(e_{i+1}\ne\bar e_i\) for every \(i\), with indices read modulo \(\ell\). It is primitive if it is not a repetition of a shorter closed walk. Let \(\mathcal P(G)\) be the set of equivalence classes of primitive cyclically nonbacktracking closed walks under cyclic rotation; reversal is not an additional equivalence. Write \(\ell(C)\) for the length of a representative. Define the Ihara zeta function as the formal power series
\[
\zeta(u,G)=\prod_{[C]\in\mathcal P(G)}\bigl(1-u^{\ell(C)}\bigr)^{-1}\in\Q[[u]].
\]
The determinant identity supplied in the record is
\[
\zeta(u,G)^{-1}=(1-u^2)^{m-n}\det\!\bigl(I_n-uA+(k-1)u^2I_n\bigr).
\]
Here the formal power series interpretation avoids imposing an unspecified analytic domain for \(u\).

\emph{Source research question.} Can techniques from the Kadison--Singer workshop, described in the supplied expanded statement as Kadison--Singer and interlacing-polynomial methods, yield further information about \(\zeta(u,G)\) for connected regular graphs?

The determinant identity is supplied background, rather than the requested new conclusion. Neither the original question nor its expanded formulation specifies a property of the zeta function to establish, a proposed bound, or the intended quantification over graphs and degrees. Thus the record defines a research direction but does not determine a single precise conjecture with a truth value. No stronger target statement is inferred here.
\subsection{Short English statement}
The question asks whether techniques discussed at the Kadison--Singer workshop, including interlacing-polynomial methods in the expanded record, can provide further information about the Ihara zeta function of a connected regular graph. The supplied determinant formula describes the object, but the desired new conclusion and some graph conventions remain unspecified. The supplied open label is a dated curation assessment.
\subsection{Sources}
\begin{itemize}
\item Supplied local UnsolvedMath registry record, permanent ID \(30002847\), title: \emph{Kadison--Singer Methods for Ihara Zeta Functions}; collection: Oberwolfach Reports Open Problems; exact source problem code: \texttt{catalogue entry}. No individual proposer is named.
\item Supplied report citation: \emph{Arbeitsgemeinschaft: The Kadison-Singer Conjecture (2016)}, Oberwolfach Reports, pp.\ 927--969, DOI \texttt{10.4171/owr/2015/17}. Supplied source URL and literature-source link, labelled ``original OWR report'': \url{https://doi.org/10.4171/owr/2015/17}. The background attributes the question to workshop open-problem discussions.
\item The registry gives proposed year \(2015\) and an \texttt{owr-year:2015} tag, while the supplied report citation includes \(2016\). This discrepancy is retained without resolution.
\item Literature assessment dated exactly \texttt{2026-08-21}: the supplied status is \texttt{open}. The supplied assessment says that the extract asks broadly for further information without specifying a proposition or target theorem, and that no definite resolution is identified. This is supplied dated curation, not an independently verified literature assessment; time-sensitive status and cited claims require independent verification.
\item The supplied statement-status label is \texttt{corrected\_verified}. Its accompanying verification note reports that the official Oberwolfach passage was checked directly and that the question, definitions, and notation were restored without adding a stronger claim. This is a supplied provenance claim; no source visit or independent verification was performed for this draft.
\end{itemize}
\end{document}
